Right-slot division makes three elementary notions meet: reciprocal
arithmetic, a pole in an affine chart, and the projective point at infinity.
Iteration adds a fourth notion: a finite or infinite expression can be the
closure of a recurrence at one of its fixed points.  This section keeps the
algebraic fixed-point equation, analytic convergence, and projective
continuation distinct.

\subsection{Arithmetic reciprocals and projective poles}

For $c\ne0$, the ordinary map
\[
  J_c(z)=\frac{c}{z}
\]
is defined on $K\setminus\{0\}$.  Its projective matrix is
\[
  J_c\leftrightarrow
  \begin{pmatrix}0&c\\1&0\end{pmatrix},
\]
so on $\PP^1(K)$ it satisfies
\[
  J_c(0)=\infty,
  \qquad
  J_c(\infty)=0.
\]
The projective statements do not make ordinary division by zero admissible;
they record the chart transition of the associated fractional linear map.

More generally, the reciprocal coordinate centered at a finite point $q$ is
\begin{equation}\label{eq:p0-pole-coordinate}
  w_q(z)=\frac1{z-q}.
\end{equation}
It sends $q$ to infinity and infinity to zero.  In this coordinate, rational
functions with no finite pole except possibly $q$ become polynomials in
$w_q$:
\[
  K[w_q]=K\!\left[\frac1{z-q}\right].
\]
The usual polynomial chart $K[z]$ is the corresponding chart whose distinguished
pole is $\infty$.  This is an elementary coordinate observation, not a claim
that one pole chart is globally preferable to another.

\begin{remark}[The pole as a ripple center]
By Theorem~\ref{thm:p0-ripple-pencil-formula}, the same point $q$ that is sent
to infinity by $w_q$ is the common boundary tangency point of the pulled-back
horocycles.  Thus the algebraic denominator zero and the geometric ripple
center are two readings of the same projective datum.
\end{remark}

\subsection{A reciprocal as the closure of an infinite expansion}

Consider a scalar parameter $r$ and the affine recursion
\begin{equation}\label{eq:p0-geometric-recursion}
  S_0=0,
  \qquad
  S_{n+1}=1+rS_n.
\end{equation}
The first terms are
\[
  0,\quad 1,\quad 1+r,\quad 1+r+r^2,\quad\ldots.
\]
As a nested right-expanded expression,
\[
  S_n=1+r\bigl(1+r(\cdots(1+r\cdot0)\cdots)\bigr).
\]
Since scalar multiplication is commutative, this example does not by itself
distinguish left and right placement; its role is to show how an arithmetic
reciprocal arises from recursive closure.

\begin{proposition}[Finite geometric expansion]
\label{prop:p0-finite-geometric-expansion}
For every $n\ge0$,
\[
  S_n=\sum_{k=0}^{n-1}r^k.
\]
If $r\ne1$, then
\begin{equation}\label{eq:p0-finite-geometric-quotient}
  S_n=\frac{1-r^n}{1-r}.
\end{equation}
If $K=\R$ or $\C$ and $|r|<1$, then
\[
  S_n\longrightarrow\frac1{1-r}.
\]
The same identity holds formally in $K[[r]]$ as
$(1-r)^{-1}=\sum_{k\ge0}r^k$.
\end{proposition}

\begin{proof}
Induction in \eqref{eq:p0-geometric-recursion} gives the finite sum.  Multiplying
that sum by $1-r$ telescopes to $1-r^n$, proving
\eqref{eq:p0-finite-geometric-quotient}.  Under $|r|<1$, one has $r^n\to0$.
The formal-power-series statement follows because multiplication by $1-r$
leaves the constant series $1$ coefficient by coefficient.
\end{proof}

The convergence hypothesis is active.  For example, $r=-1$ makes the sequence
oscillate, and $|r|>1$ generally makes its finite values grow rather than
converge.  The algebraic reciprocal $1/(1-r)$ may exist even when the forward
iteration does not converge to it.

\subsection{Fixed-point closure and escape to infinity}

Let
\[
  T_r(s)=1+rs.
\]
A finite fixed point satisfies
\[
  s=T_r(s)
  \quad\Longleftrightarrow\quad
  (1-r)s=1.
\]
Hence
\begin{equation}\label{eq:p0-geometric-fixed-point}
  q_r=\frac1{1-r}
\end{equation}
when $r\ne1$.  The geometric series therefore converges, when $|r|<1$, to the
finite fixed point of the update that generates its truncations.

The projective matrix of $T_r$ is
\[
  M_r=
  \begin{pmatrix}r&1\\0&1\end{pmatrix}.
\]
It always fixes $\infty$.  If $r\ne1$, it has the second fixed point $q_r$.
As $r\to1$, the finite point $q_r$ leaves every compact subset of the affine
line and approaches $\infty$ in $\PP^1$.  At $r=1$ the update becomes the
translation
\[
  T_1(s)=s+1,
\]
which has no finite fixed point and has $\infty$ as its repeated projective
fixed point.  Thus the pole of the parameter function $r\mapsto1/(1-r)$ is
exactly the parameter at which finite fixed-point closure degenerates into a
parabolic fixed point at infinity.

The same mechanism holds for a general invertible affine map.

\begin{proposition}[Affine iteration and its fixed point]
\label{prop:p0-affine-iteration}
Let
\[
  T(z)=az+b,
  \qquad a\ne0.
\]
If $a\ne1$, put
\[
  q=\frac{b}{1-a}.
\]
Then
\begin{equation}\label{eq:p0-affine-centered-form}
  T(z)=q+a(z-q),
  \qquad
  T^n(z)=q+a^n(z-q).
\end{equation}
Consequently, over $\R$ or $\C$, $T^n(z)\to q$ for every finite $z$ when
$|a|<1$.  If $a=1$ and $b\ne0$, the only projective fixed point is
$\infty$ and $T^n(z)=z+nb$.
\end{proposition}

\begin{proof}
The equation $(1-a)q=b$ gives the centered form.  Iteration of that form gives
\eqref{eq:p0-affine-centered-form} by induction.  The convergence and
translation statements are immediate.
\end{proof}

This proposition separates three notions:

\begin{itemize}
  \item the algebraic existence of a finite fixed point;
  \item attraction of an orbit to that point;
  \item the projective fixed point at infinity already carried by every affine
  operator.
\end{itemize}

They coincide only under additional hypotheses.

\subsection{Right reciprocal iteration and continued fractions}

A stationary right-expanded reciprocal update has the form
\begin{equation}\label{eq:p0-right-reciprocal-map}
  F(z)=a+\frac{b}{z},
  \qquad b\ne0,
\end{equation}
with projective matrix
\[
  M_F=
  \begin{pmatrix}a&b\\1&0\end{pmatrix}.
\]
The ordinary iteration is defined only while no intermediate denominator
vanishes.  Its finite truncations are generalized continued fractions:
\[
  z_n
  =a+\cfrac{b}{a+\cfrac{b}{\ddots+\cfrac{b}{z_0}}}.
\]
The matrix formulation remains projectively meaningful when an affine
intermediate value passes through infinity.

\begin{proposition}[Fixed points of a reciprocal update]
\label{prop:p0-reciprocal-fixed-points}
The finite fixed points of \eqref{eq:p0-right-reciprocal-map} are the roots of
\begin{equation}\label{eq:p0-reciprocal-quadratic}
  z^2-az-b=0.
\end{equation}
At a nonzero fixed point $q$,
\[
  F'(q)=-\frac{b}{q^2}.
\]
Over $\C$, if $|F'(q)|<1$, then $q$ is locally attracting; if
$|F'(q)|>1$, it is locally repelling.
\end{proposition}

\begin{proof}
Multiplying $q=a+b/q$ by $q$ gives
\eqref{eq:p0-reciprocal-quadratic}.  Differentiation gives the multiplier.
The local dynamical classification is the standard contraction/expansion test
from the first nonzero derivative.
\end{proof}

\begin{example}[The golden-ratio expansion]
\label{ex:p0-golden-ratio}
Take
\[
  F(z)=1+\frac1z,
  \qquad z_0=1.
\]
Then
\[
  1,\quad2,\quad\frac32,\quad\frac53,\quad\frac85,\ldots
\]
are ratios of consecutive Fibonacci numbers.  The fixed points solve
\[
  z^2-z-1=0,
\]
so they are
\[
  \phi=\frac{1+\sqrt5}{2},
  \qquad
  -\phi^{-1}=\frac{1-\sqrt5}{2}.
\]
Since $|F'(\phi)|=\phi^{-2}<1$, the positive fixed point is attracting for
nearby positive initial values.  The infinite continued fraction
\[
  1+\cfrac1{1+\cfrac1{1+\cdots}}
\]
therefore denotes $\phi$ only after the convergence statement has been
established; the fixed-point equation alone does not choose the attracting
root in every field or domain.
\end{example}

For nonstationary data,
\[
  F_k(z)=a_k+\frac{b_k}{z},
  \qquad
  M_k=\begin{pmatrix}a_k&b_k\\1&0\end{pmatrix},
\]
and the $n$th truncation is represented by the product
\[
  M_nM_{n-1}\cdots M_1.
\]
This is the first place where the nested visual form, the ordinary recurrence,
and matrix multiplication become exactly the same calculation.

\subsection{Fixed points of a general projective operator}

Let
\[
  F(z)=\frac{Az+B}{Cz+D},
  \qquad \Delta=AD-BC\ne0.
\]
A finite fixed point satisfies
\begin{equation}\label{eq:p0-general-mobius-fixed-point}
  Cz^2+(D-A)z-B=0.
\end{equation}
In homogeneous coordinates, fixed points are precisely the eigenlines of the
matrix
\[
  \begin{pmatrix}A&B\\C&D\end{pmatrix}.
\]
The point $\infty$ is fixed exactly when $C=0$.  Over an algebraic closure a
nonidentity projective map has two fixed points counted with multiplicity; over
the base field those points may be distinct, repeated, or visible only after a
field extension.

This completes the elementary circle of ideas:
\[
  \text{reciprocal}
  \longrightarrow
  \text{pole}
  \longrightarrow
  \infty
  \longrightarrow
  \text{fixed-point closure}
  \longrightarrow
  \text{matrix eigenline}.
\]
The next section shows that the same matrix also governs point paths and ripple
pencils.
